\documentclass[11pt]{article}
\usepackage[T1]{fontenc}
\usepackage{textcomp}
\usepackage[margin=0.75in]{geometry}
\usepackage{amsmath,amssymb,amsthm}
\usepackage{array,booktabs}
\usepackage{hyperref}
\setlength{\parindent}{0pt}
\newtheorem{theorem}{Theorem}
\theoremstyle{definition}
\newtheorem{definition}{Definition}
\title{Flow Proof Artifact: prop-21-duplicate-ratio-homologous-sides}
\author{Generated by \texttt{flow doc proof}}
\date{Flow Proof Book}
\begin{document}
\maketitle
Similar rectilineal figures are to one another as the duplicate ratio of their homologous sides.\\[0.5em]
\textbf{Source.} euclid — Elements, Book VI, Proposition 21\\[0.5em]
\bigskip
\noindent{\large \textbf{Derived fact 1.} \textit{Similar rectilineal figures are to one another as the duplicate ratio of their homologous sides} \hfill the Euclidean plane \cdot Euclid Book VI \cdot Proposition 21: similar figures have ratio duplicate of homologous sides}\par
\smallskip\noindent\textit{Coordinate: the Euclidean plane \cdot Euclid Book VI \cdot Proposition 21: similar figures have ratio duplicate of homologous sides}\par
\smallskip\noindent\textit{Source: euclid — Elements, Book VI, Proposition 21}\par
\smallskip\noindent\textit{Built on: proposition 20: similar polygons are divided into the same number of similar triangles, for Euclid Book VI on the Euclidean plane}\par
\medskip
\noindent\textbf{Goal.} Similar rectilineal figures are to one another as the duplicate ratio of their homologous sides\par
\noindent$\displaystyle \text{area as duplicate homologous ratio}$\par
\medskip
\noindent
\begin{tabular}{@{} >{\raggedright\arraybackslash}p{0.06\textwidth} >{\raggedright\arraybackslash}p{0.47\textwidth} >{\raggedleft\arraybackslash}p{0.41\textwidth} @{}}
\textbf{\#} & \textbf{Proof} & \textbf{Mathematics} \\
\midrule
\textbf{1.} & We prove that proposition 21: similar figures have ratio duplicate of homologous sides for Euclid Book VI the Euclidean plane. & \\[0.45em]
\textbf{2.} & Let fig1 = figure\_1. & \\[0.45em]
\textbf{3.} & Let fig2 = figure\_2. & \\[0.45em]
\textbf{4.} & Let s1 = homologous\_side\_1. & \\[0.45em]
\textbf{5.} & From step 2, step 3, and step 4, we can deduce that ratio areas equals duplicate ratio(s1, s2). & $\displaystyle \text{ratio areas} = \text{duplicate ratio}(s1, s2)$ \\[0.45em]
\textbf{6.} & We invoke the derived fact governing Euclid Book VI the Euclidean plane: proposition 20: similar polygons are divided into the same number of similar triangles, for Euclid Book VI on the Euclidean plane. & \\[0.45em]
\textbf{7.} & From step 6, this implies area as duplicate homologous ratio. Hence proven. & $\displaystyle \text{area as duplicate homologous ratio}$ \\[0.45em]
\bottomrule
\end{tabular}
\label{thm:Geometry-Euclid Book VI-proposition_21_similar_figures_have_ratio_duplicate_of_homologous_sides}
\par\medskip\hrule
\end{document}
